\section{Experiments and Results}
\label{sec:experiments}
\subsection{Evaluation Questions and Setup}
State which claim each experiment tests. Describe the data or study
materials, inclusion criteria, baselines, evaluation metrics, and procedure.
Explain how comparisons are made fair and how uncertainty is assessed.
Report actual experimental conditions instead of inventing missing details.

\subsection{Quantitative Evaluation}
Table~\ref{tab:results} demonstrates a compact results table.
All entries are intentionally unfilled; dashes do not mean zero.
Replace column headings with task-specific metrics and units, and indicate
whether larger or smaller values are preferred.

\begin{table}[t]
  \centering
  \caption{Placeholder results. All dashes mean ``not measured.''
  Replace method names, metric headings, and entries with actual results.}
  \label{tab:results}
  \begin{tabular}{lcc}
    \toprule
    Method & Metric A & Metric B \\
    \midrule
    Baseline A & --- & --- \\
    Baseline B & --- & --- \\
    Your method & --- & --- \\
    \bottomrule
  \end{tabular}
\end{table}


Describe the main observation supported by your measurements. Separate
observations from explanations, and avoid claiming an improvement without
appropriate evidence. If experiments have not been conducted, label this
section as an evaluation plan and do not present anticipated outcomes as results.

\subsection{Qualitative Evaluation}
Use examples to explain behavior that a single numerical metric cannot
capture. Figure~\ref{fig:qualitative} provides a two-panel placeholder.
Include informative captions, matched conditions, and both successful and
unsuccessful cases where relevant.

\begin{figure}[t]
  \centering
  \fbox{\begin{minipage}[c][26mm][c]{0.43\columnwidth}
    \centering\textbf{DUMMY A}\\[3mm]Input / baseline
  \end{minipage}}\hfill
  \fbox{\begin{minipage}[c][26mm][c]{0.43\columnwidth}
    \centering\textbf{DUMMY B}\\[3mm]Output / comparison
  \end{minipage}}
  \caption{Placeholder qualitative comparison. Replace both panels with
  your own examples and explain the relevant differences. These empty
  panels do not represent experimental evidence.}
  \label{fig:qualitative}
\end{figure}


\subsection{Ablation or Sensitivity Analysis}
Test the role of a design choice by changing one factor while holding other
conditions fixed. Explain what the comparison reveals and what it cannot
establish. If this analysis is inappropriate for your study, replace it with
a suitable robustness check or alternative analysis.
